% refs/nnd_mcmc.tex
% Salient algorithmic details from Segert & Wycoff (2025)
% "A Probabilistic Basis for Low-Rank Matrix Learning"  arXiv:2510.05447
% Adapted to the heteroscedastic denoising/completion setting in matlap.

\documentclass[11pt]{article}

\usepackage[utf8]{inputenc}
\usepackage[T1]{fontenc}
\usepackage{mathpazo} % Palatino font
\usepackage{amsmath, amssymb}
\usepackage{geometry}
\geometry{a4paper, margin=1.2in}

\title{MCMC Samplers for the Nuclear-Norm Distribution Posterior\\
\large Notes on Segert \& Wycoff (2025), arXiv:2510.05447}
\author{}
\date{}

\begin{document}
\maketitle

\section*{MCMC Samplers for the Nuclear-Norm Distribution Posterior}

\subsection*{Source paper}
Segert, J.\ \& Wycoff, N.\ (2025).
\emph{A Probabilistic Basis for Low-Rank Matrix Learning}.
arXiv:2510.05447v1.

\subsection*{Model}
Observation model (heteroscedastic; entries missing when $s_{ij}=\infty$):
\[
  y_{ij} \mid x_{ij},\, s_{ij}
  \;\sim\; \mathcal{N}(x_{ij},\, s_{ij}^2), \quad (i,j)\in\Omega.
\]
Nuclear-norm prior with scalar $\lambda>0$ (Proposition~3.3 of source):
\[
  p(X \mid \lambda) \;\propto\; \lambda^{mn}\exp\!\bigl(-\lambda\|X\|_*\bigr).
\]
Half-Cauchy hyperprior on $\lambda$ (Section~4.3 of source; Polson \& Scott 2012):
implemented via the inverse-Gamma hierarchy of Makalic \& Schmidt (2015)---see
Section~\ref{sec:lambda} below.

\subsection*{SVD reparameterization}
Write the thin SVD $X = U\operatorname{diag}(\sigma)V^T$ where
$U\in\mathrm{St}(m,n)$ (semi-orthogonal), $V\in O(n)$ (orthogonal),
$\sigma\in\mathbb{R}^n_+$.  The nuclear norm is $\|X\|_*=\sum_k\sigma_k$.
Reparameterizing to $(U,\sigma,V)$ gives tractable conditional posteriors
(Section~4.2 of source; homoscedastic $s_{ij}^2=\gamma^2$).


% ===========================================================================
\section*{Sampler 1: Proximal MALA (Sec.\ 4.1, Appendix J.4)}
% ===========================================================================

\paragraph{Algorithm.}
Metropolis-Adjusted Langevin Algorithm using the proximal operator of
$-\log p$ as the proposal mean.  For the heteroscedastic setting, the
proximal operator of $\lambda\|\cdot\|_*$ does not simplify to a closed-form
expression (unlike the homoscedastic case in Appendix~J.3 where the proximal
operator of the negative log-posterior has a closed form via SVT).  We
therefore use a \emph{proximal gradient} step instead (as in the matrix
completion case, Appendix~J.4):

\[
  \bar X = \mathrm{SVT}_{\delta\lambda/2}\!\Bigl(
    X^t + \tfrac{\delta}{2}\, \underbrace{M\odot(Y-X^t)/S^2}_{
      \nabla_X\log p_{\text{lik}}(X^t)}
  \Bigr),
\]
where $M_{ij}=\mathbf{1}[(i,j)\in\Omega]$, $\mathrm{SVT}_\tau(Z)=
U\operatorname{diag}(\max(\sigma_k-\tau,0))V^T$ is singular-value
soft-thresholding, and $\delta>0$ is the proposal scale.

\paragraph{Proposal and accept/reject.}
\[
  X^* \;\sim\; \mathcal{N}(\bar X,\, \delta I),
\]
accepted with probability $\min(1,\alpha)$ where
\[
  \log\alpha
  = \log p(X^*) - \log p(X^t)
  + \log q(X^t\mid X^*) - \log q(X^*\mid X^t),
\]
\[
  \log p(Z)
  = -\tfrac{1}{2}\sum_{(i,j)\in\Omega}\frac{(y_{ij}-z_{ij})^2}{s_{ij}^2}
  - \lambda\|Z\|_*,
\]
\[
  \log q(Z'\mid Z) = -\frac{\|Z'-\bar Z\|_F^2}{2\delta},
\]
with $\bar Z$ the proximal-gradient mean centred at $Z$.

\paragraph{Step size adaptation.}
Robbins--Monro targeting acceptance rate $0.574$ (optimal for proximal MCMC;
Crucinio et al.\ 2023):
\[
  \log\delta_{t+1} = \log\delta_t + c_t(\alpha_t - 0.574),
  \quad c_t = c/\sqrt{t}.
\]

\paragraph{Heteroscedastic adaptation.}
In the homoscedastic case (Appendix~J.3), the proximal operator of the full
negative log-posterior has the closed form
$\mathrm{SVT}_{\delta\lambda(\delta+2\gamma^2)^{-1}}
 \bigl((\delta Y + 2\gamma^2 X^t)(\delta+2\gamma^2)^{-1}\bigr)$.
For heteroscedastic $s_{ij}^2$, no such closed form exists; the proximal
gradient approximation above is used instead, exactly as in the matrix
completion variant of Appendix~J.4.  The MH correction renders the chain
exact regardless of this approximation.


% ===========================================================================
\section*{Sampler 2: SVD Gibbs (Sec.\ 4.2)}
\label{sec:gibbs}
% ===========================================================================

\paragraph{Algorithm.}
Iterate over three conditional samplers in the SVD parameterization
$X=U\operatorname{diag}(\sigma)V^T$.  Let $s=\operatorname{diag}(U^TYV)$.

\subparagraph{Step A: Sample $\sigma\mid U,V,\lambda,Y$.}
The joint conditional density of the singular values is
\[
  p(\sigma\mid\cdot)
  \;\propto\;
  \exp\!\Bigl(-\lambda\textstyle\sum_k\sigma_k
    -\tfrac{\|\sigma - s\|^2}{2\gamma^2}\Bigr)
  \prod_{k\le n}\sigma_k^{m-n}
  \prod_{i<j\le n}|\sigma_i^2 - \sigma_j^2|,
\]
where the Vandermonde factor $\prod_{i<j}|\sigma_i^2-\sigma_j^2|$ arises from
the Jacobian of the SVD map.  This density is not product-separable and
has posterior nonsmoothness from the nuclear norm; the source samples it via
a Metropolis-adjusted Langevin method operating in $\mathbb{R}^n$ (vs.\ the
full $m\times n$ problem for direct MALA on $X$).

\subparagraph{Step B: Sample $U\mid\sigma,V,\lambda,Y$.}
\[
  p(U\mid\cdot)
  \;\propto\;
  \exp\!\Bigl(-\tfrac{1}{2\gamma^2}\operatorname{tr}\bigl[(YV\operatorname{diag}(\sigma))^T U\bigr]\Bigr),
\]
which is a \emph{matrix von Mises--Fisher} (vMF) distribution on
$\mathrm{St}(m,n)$ (Jupp \& Mardia 1979).  Sampled via the reflective slice
sampler of Hoff (2009).

\subparagraph{Step C: Sample $V\mid\sigma,U,\lambda,Y$.}
\[
  p(V\mid\cdot)
  \;\propto\;
  \exp\!\Bigl(-\tfrac{1}{2\gamma^2}\operatorname{tr}\bigl[(\operatorname{diag}(\sigma)Y^TU^T)^T V\bigr]\Bigr),
\]
again a matrix vMF on $O(n)$; sampled via Hoff (2009).

\paragraph{The variational bound is not a scale mixture.}
\label{sec:het-gibbs}
Proposition~3.2 of the source gives a variational (minimax) identity:
\[
  \exp(-\lambda\|X\|_*)
  \;=\;
  \min_{Q\succ 0}
  \exp\!\Bigl(-\tfrac{\lambda}{2}\bigl[\operatorname{tr}(Q^{-1}\Psi)
  + \operatorname{tr}(Q)\bigr]\Bigr),
  \quad \Psi = X^TX,
\]
minimized at $Q_* = \Psi^{1/2}$ by the AM--GM inequality
$\operatorname{tr}(Q^{-1}\Psi)+\operatorname{tr}(Q)\ge 2\operatorname{tr}(\Psi^{1/2})=2\|X\|_*$.
This is a \emph{minimization}, not a marginalization: the
\emph{integral} $\int_{Q\succ 0}\exp\bigl(-\tfrac{\lambda}{2}
[\operatorname{tr}(Q^{-1}\Psi)+\operatorname{tr}(Q)]\bigr)dQ$
does \emph{not} equal $\exp(-\lambda\|X\|_*)$ in general.
A Gibbs chain that treats $Q$ as an auxiliary variable and samples from
the ``pseudo-joint''
$p(X,Q)\propto\exp(-\tfrac{\lambda}{2}[\operatorname{tr}(Q^{-1}\Psi)+\operatorname{tr}(Q)])\cdot
p(Y|X)$
therefore targets the \emph{wrong} posterior.

\paragraph{Correct Gaussian scale mixture.}
\label{sec:gsm-correct}
The proper augmentation uses the modified Bessel identity
$K_{1/2}(z)=\sqrt{\pi/(2z)}\,e^{-z}$:
\[
  \int_0^\infty d^{-1/2}
  \exp\!\Bigl(-\tfrac{\psi}{2d} - \tfrac{\lambda^2 d}{2}\Bigr)\,dd
  \;=\; \frac{\sqrt{\pi}}{\lambda}\,e^{-\lambda\sqrt{\psi}}
  \;=\; \frac{\sqrt{\pi}}{\lambda}\,e^{-\lambda\sigma},
  \quad \sigma=\sqrt{\psi},
\]
so that
\[
  \exp(-\lambda\|X\|_*)
  \;=\; \exp\!\Bigl(-\lambda\textstyle\sum_k\sigma_k\Bigr)
  \;\propto\;
  \int_{\mathbb{R}^n_{>0}}
  \prod_k d_k^{-1/2}
  \exp\!\Bigl(-\tfrac{\psi_k}{2d_k} - \tfrac{\lambda^2 d_k}{2}\Bigr)
  \prod_k dd_k,
\]
where $\psi_k$ are the eigenvalues of $\Psi=X^TX$ (i.e.\ $\psi_k=\sigma_k^2$).
Crucially, $\lambda$ appears \emph{quadratically} in the exponential
(not linearly as in the variational bound).  The auxiliary variable $Q$ must
share eigenvectors with $\Psi$ -- write $Q=V\operatorname{diag}(d)V^T$ where
$\Psi=V\operatorname{diag}(\psi)V^T$ -- for the integral to factorise.
The correct conditional is
\[
  \boxed{p(d_k\mid X,\lambda)
  \;\propto\; d_k^{-1/2}\exp\!\Bigl(-\tfrac{\psi_k}{2d_k}-\tfrac{\lambda^2 d_k}{2}\Bigr)
  \;=\;\mathrm{GIG}\!\Bigl(\tfrac{1}{2},\,\psi_k,\,\lambda^2\Bigr).}
\]
This differs from the variational-bound ``conditional'' in two ways:
the shape parameter is $\tfrac{1}{2}$ (giving the $d_k^{-1/2}$ factor) and the
rate on $d_k$ is $\lambda^2/2$, not $\lambda/2$.

\paragraph{Heteroscedastic Gibbs (corrected).}
For heteroscedastic $s_{ij}^2$, the vMF structure of the conditionals for
$U$ and $V$ is lost.  Using the correct GSM with auxiliary $d$ in the
eigenbasis of $X^TX$, the three-step Gibbs is:

\subparagraph{Step A: Sample $X\mid d,\lambda,Y,S$.}
Writing $p(x_i\mid Q,\lambda)\propto\exp(-\tfrac{1}{2}x_i^TQ^{-1}x_i)$
with precision $\lambda Q^{-1}$, each row is a Gaussian:
\begin{align}
  A_i &= \operatorname{diag}(s_{i,\cdot}^{-2}\odot m_{i,\cdot}) + \lambda Q^{-1},
  \quad
  r_i = \bigl(s_{i,\cdot}^{-2}\odot m_{i,\cdot}\bigr)\odot y_i^{\mathrm{obs}},
  \\
  x_i &\;\sim\; \mathcal{N}(A_i^{-1}r_i,\; A_i^{-1}),
\end{align}
where $Q^{-1}=V\operatorname{diag}(1/d)V^T$.
Solved via Cholesky of the $n\times n$ precision $A_i$; rows parallelised
via \texttt{vmap}.

\subparagraph{Step B: Sample $d\mid X,\lambda$ (exact GIG).}
After sampling $X_\mathrm{new}$, eigendecompose
$\Psi_\mathrm{new}=V_\mathrm{new}\operatorname{diag}(\psi_\mathrm{new})V_\mathrm{new}^T$.
Each $d_k$ is independently drawn from
\[
  d_k\mid\psi_k,\lambda
  \;\sim\; \mathrm{GIG}\!\Bigl(\tfrac{1}{2},\,\psi_k,\,\lambda^2\Bigr),
\]
with density $\propto d_k^{-1/2}\exp(-\psi_k/(2d_k)-\lambda^2 d_k/2)$.
GIG samples can be drawn exactly (no Metropolis step needed) using the
algorithm of Devroye (2014) or Hörmann \& Leydold (2014); see Appendix.
The mode is $d_k^* = (-\tfrac{1}{2}+\sqrt{\tfrac{1}{4}+\lambda^2\psi_k})/\lambda^2$
(close to $\sigma_k/\lambda$ for large $\lambda\sigma_k$).

\subparagraph{Step C: Sample $\lambda\mid d_\mathrm{new}$ (half-Cauchy hierarchy).}
From the correct GSM joint,
$p(\lambda\mid d)\propto\lambda^{(m+1)n}\exp(-\tfrac{\lambda^2}{2}\textstyle\sum_k d_k)\cdot p(\lambda)$,
where the $\lambda^{(m+1)n}$ factor accounts for the $\lambda^{mn}$ normalising
constant of the NND prior and the $\lambda^n$ from the GIG normalising constants.
Adopting a half-Cauchy prior on $\lambda$ via the half-Normal scale mixture
(Polson \& Scott 2012): $\lambda|\zeta\sim\mathrm{half\text{-}Normal}(0,\zeta)$,
$\zeta\sim\mathrm{IG}(\tfrac{1}{2},\tfrac{1}{2})$:
\begin{align}
  \zeta' &\;\sim\; \mathrm{IG}\!\Bigl(\tfrac{1}{2},\;\tfrac{1+\lambda^2}{2}\Bigr),
  \label{eq:zeta-update}\\
  \lambda' &\;\sim\; \mathrm{TruncatedNormal}^+\!\Bigl(
    0,\;\bigl[\textstyle\sum_k d_k' + 1/\zeta'\bigr]^{-1}\Bigr),
  \label{eq:lambda-update}
\end{align}
where TruncatedNormal$^+(0,v)$ denotes a half-normal with variance $v$,
sampled as $|\mathcal{N}(0,1)|/\sqrt{\sum d_k'+1/\zeta'}$.
The $(m+1)n$ exponent modifies \eqref{eq:lambda-update} to a
truncated generalised normal; if $(m+1)n$ is large this can be approximated
or handled via MH.

\paragraph{Comparison with current \textsc{matlap} implementation.}
The current code uses the variational pseudo-joint rather than the correct
scale mixture, giving two errors in Step~B:
(i) the GIG shape is $1$ instead of $\tfrac{1}{2}$ (missing the $d_k^{-1/2}$
factor), and (ii) the rate on $d_k$ is $\lambda/2$ instead of $\lambda^2/2$.
Step C also uses the $\exp(-\lambda\,\operatorname{tr}(Q))$ form (linear in $\lambda$)
rather than the quadratic-in-$\lambda$ form from the correct scale mixture.
The sampler is therefore approximate; the stationary distribution differs
from the true posterior.

\paragraph{Sparse-data behaviour.}
At high missing fractions ($\gtrsim$60\%), $X_\mathrm{new}\to 0$ under the
posterior, so $\psi_k\to 0$, $\sum d_k\to 0$ and $\lambda'\to\infty$.
This occurs with both the variational and the correct GSM formulations:
the NND prior places all mass at $X=0$ for large $\lambda$, and sparse
data cannot rule out this degenerate configuration.

% ===========================================================================
\section*{$\lambda$ sampling: Half-Cauchy prior (Sec.\ 4.3)}
\label{sec:lambda}
% ===========================================================================

A Gamma prior on $\lambda$ is conjugate but may not allow $\lambda$ to range
over multiple orders of magnitude.  The source uses a \emph{half-Cauchy}
prior (Polson \& Scott 2012), implemented via the inverse-Gamma hierarchy of
Makalic \& Schmidt (2015): if $\alpha|\beta\sim\mathrm{IG}(1/2, 1/\beta)$ and
$\beta\sim\mathrm{IG}(1/2,1)$ then $\alpha$ is marginally Half-Cauchy.  For
conjugacy with the NND likelihood $p(X|\lambda)\propto\lambda^{mn}\exp(-\lambda\|X\|_*)$,
the source samples the \emph{inverse} $\alpha=1/\lambda$ and introduces an
auxiliary $\beta$.  In the SVD Gibbs (homoscedastic), the conjugate updates are:
\begin{align}
  \beta|\alpha &\;\sim\; \mathrm{IG}\!\Bigl(1,\; 1 + \tfrac{1}{\alpha}\Bigr), \\
  \alpha|\beta,X,\gamma^2 &\;\sim\;
  \mathrm{IG}\!\Bigl(mn + \tfrac{1}{2},\;
    \tfrac{\|X\|_*}{\gamma^2} + \tfrac{1}{\beta}\Bigr), \\
  \lambda &\;\leftarrow\; \tfrac{1}{\alpha}.
\end{align}

\paragraph{Note on parameterisations.}
The IG hierarchy above is conjugate to the \emph{linear} NND likelihood
$\exp(-\lambda\|X\|_*/\gamma^2)$ (source uses homoscedastic $\gamma^2$).
In the correct GSM augmentation (Section~\ref{sec:gibbs}), $\lambda$
appears \emph{quadratically} via $\exp(-\lambda^2\sum_k d_k/2)$,
so the conditional $p(\lambda|d)$ has a different form (see Step~C in
Section~\ref{sec:gibbs}) and the IG hierarchy no longer applies.
The current \textsc{matlap} implementation applies the IG hierarchy
with $\operatorname{tr}(Q) = \sum_k d_k$ in place of $\|X\|_*/\gamma^2$;
this is internally consistent with the (approximate) variational pseudo-joint
but not with the correct GSM.

\end{document}
